% SPDX-License-Identifier: Apache-2.0
\section{Units, Processes and Parallelism}
\label{sec:e-units}

A unit is a struct that implements \code{crate::comp::Module}. The trait
has one method.

\begin{lstlisting}[caption={The unit trait. Inputs and outputs are type parameters, so one struct may implement it more than once.},label={lst:module}]
#[allow(async_fn_in_trait)]
pub trait Module<In, Out> {
    async fn run(&mut self, inputs: In, outputs: Out);
}
\end{lstlisting}

Inputs and outputs are type parameters rather than associated types, and
that choice buys something concrete: one struct may implement the trait
more than once, with different shapes, so a unit that works on 32-bit and
64-bit streams is one implementation and not two.

\subsection{Parallelism is stated, not implied}

Synthesis calls \code{run}. A unit with more than one concurrent process
starts one async function per process and joins them.

\begin{lstlisting}[caption={Parallelism written down. A unit with one process writes no join.},label={lst:join}]
impl Module<Ins, Outs> for DualPort {
    async fn run(&mut self, _i: Ins, _o: Outs) {
        join2(self.port_a(), self.port_b()).await;
    }
}
\end{lstlisting}

This replaces a rule with a construct. TxHDL said that several sequence
blocks in one module run concurrently, which was a property of having
declared them. It also contradicted itself: the specification stated that a
module holds one sequence, while three of its own modules declared
two~\cite{unification}. Here there is no rule to contradict. Parallelism is
a \code{join}, written where it happens, and a unit with one process does
not write one.

\subsection{Processes cannot take a mutable borrow}

The obvious way to write the two processes does not compile, and finding
that out cost ten lines.

\begin{lstlisting}[caption={What does not compile, and the error it gives. This is the code a Rust programmer writes first.},label={lst:e0499}]
async fn port_a(&mut self) { /* ... */ }
async fn port_b(&mut self) { /* ... */ }

// error[E0499]: cannot borrow `*self` as mutable
//               more than once at a time
join2(self.port_a(), self.port_b()).await;
\end{lstlisting}

Hardware processes share state: they drive registers. Rust permits one
mutable borrow at a time. The two facts collide directly.

The repair is that a register is a cell, and a process takes a shared
reference.

\begin{lstlisting}[caption={The repair. A register is a cell several processes reach, which is what a register is. Reading one is \code{async}; driving one is not.},label={lst:reg}]
pub struct Reg<T: Copy>(Cell<T>);

impl<T: Copy> Reg<T> {
    pub async fn get(&self) -> T { tick().await; self.0.get() }
    pub fn set(&self, x: T) { self.0.set(x) }
    pub fn peek(&self) -> T { self.0.get() }   // testbench only
}

async fn port_a(&self) {
    loop {
        let hits = self.hits_a.get().await;
        self.hits_a.set(hits + 1);
    }
}
\end{lstlisting}

\code{run} still takes \code{\&mut self}, and the processes it starts take
\code{\&self}. That compiles, and it also describes the hardware better: a
register is a thing several processes reach, which is what a register is.

\subsection{A register read is the wait, and a process loops}
\label{sec:e-regread}

Listing~\ref{lst:reg} makes two claims that earlier drafts did not, and
both follow from what a register is.

\textbf{Reading a register is \code{async}.} A wire has no clock edge, so
\code{In::get} is plain and returns at once. A register holds the value
latched at the last edge, so reading it is waiting for that edge, and
\code{Reg::get} is \code{async} to say so. Driving a register states the
next value and waits for nothing, so \code{set} is plain. The
\code{.await} on a read is the mark the lowering turns into the
register; the value alive across it is what gets stored. \code{peek}
reads without waiting, and exists for testbenches, for the reason
\code{assert!} exists: it observes and drives nothing, and synthesis
skips it.

\textbf{A process loops.} A unit runs for as long as the clock does, so
every leaf \code{run} is a \code{loop}, and each iteration contains a
register read. That read is where the iteration waits for the edge, so
one iteration is one cycle. A parent whose \code{run} only joins its
children needs no loop of its own: the children loop, the join never
completes, and the parent is continuous through them.

The prototype enforces this in the only way a simulation can. A register
read yields to the executor once, so one poll of the top unit is one
cycle, and a loop that reads no register and awaits no operator never
yields and spins. That is the right behaviour, because such a loop is a
combinational loop in the hardware, which is the thing that was wrong.

\begin{pitfall}
This is the constraint most likely to surprise somebody writing a unit for
the first time, because the code that fails is the code a Rust programmer
would write first. It should be in the library's documentation and in the
error message, not only in a paper.
\end{pitfall}


\subsection{A unit that contains units}
\label{sec:e-hierarchy}

A submodule is a field. A bus that joins two submodules is a field too,
owned by whichever unit contains both of its ends.

\begin{lstlisting}[caption={Submodules and the wire between them are fields of the parent. Each child is generic over its own configuration, not the parent's.},label={lst:hier}]
pub struct Top<TC: TopConfig> {
    pub producer: Producer<TC::P>,
    pub consumer: Consumer<TC::C>,
    pub link: DataBus,
    pub pes: [Pe; 4],
}
\end{lstlisting}

\subsection{Configurations nest}
\label{sec:e-nestedconfig}

Section~\ref{sec:e-binding} gives a design one configuration and makes
every unit generic over it. That holds while every choice is global. It
stops holding as soon as a submodule has choices of its own, and
Listing~\ref{lst:hier} is where that shows.

A configuration is a tree. A unit declares the configuration \emph{it}
needs, and a parent's configuration names its children's rather than
holding their fields.

\begin{lstlisting}[caption={Each unit declares what it needs. The parent names its children's configurations and never mentions their knobs.},label={lst:nestedcfg}]
pub trait ProducerConfig {
    type Gen: Generator;
    const DEPTH: usize;
}

pub trait ConsumerConfig {
    const DEPTH: usize;        // the same name, a different unit
    const CHECKSUM: bool;
}

pub trait TopConfig {
    type P: ProducerConfig;
    type C: ConsumerConfig;
    const NAME: &'static str;
}
\end{lstlisting}

Two things follow, and both were the reason to do it.

\textbf{Two children may want a knob of the same name.} Both
\code{DEPTH}s above exist, differ, and never meet, because each belongs
to its own trait. A flat configuration would have had to rename one of
them after the fact, which is a rename in a design that had nothing wrong
with it.

\textbf{A child's configuration is reusable.} It is written once and
named by whichever parents want it, so two builds may differ in one
subsystem and share the rest.

\begin{lstlisting}[caption={Two builds. The consumer's configuration is written once and used by both; only the producer differs.},label={lst:nestedbuilds}]
pub struct CheckedConsumer;
impl ConsumerConfig for CheckedConsumer {
    const DEPTH: usize = 64;
    const CHECKSUM: bool = true;
}

pub struct Fpga;
impl TopConfig for Fpga {
    type P = FastProducer;
    type C = CheckedConsumer;
    const NAME: &'static str = "fpga";
}

pub struct Tiny;
impl TopConfig for Tiny {
    type P = SmallProducer;
    type C = CheckedConsumer;   // reused, unchanged
    const NAME: &'static str = "tiny";
}
\end{lstlisting}

Nesting costs nothing at elaboration. A child's constant is still a
compile-time value reached through the parent, so it may still size an
array:

\begin{lstlisting}[caption={A child's constant, reached through the parent, still sizes an array. This is the long spelling.},label={lst:nestedconst}]
pub const DEPTH: usize =
    <<Fpga as TopConfig>::P as ProducerConfig>::DEPTH;
pub static BUF: [u32; DEPTH] = [0; DEPTH];
\end{lstlisting}

The spelling in Listing~\ref{lst:nestedconst} is the one real cost of
nesting: reaching two levels down needs a fully qualified path, and three
levels reads worse again. The repair is one type alias per level of the
tree, written once beside the trait, so it serves every build rather than
one.

\begin{lstlisting}[caption={The repair. A generic alias per level, and a per-build alias where one build is being talked about. Both compile in probe~17.},label={lst:nestedalias}]
pub type ProducerOf<T> = <T as TopConfig>::P;
pub type ConsumerOf<T> = <T as TopConfig>::C;

// One qualification, not two, for any build.
pub const DEPTH: usize =
    <ProducerOf<Fpga> as ProducerConfig>::DEPTH;

// None at all, where one build is meant.
pub type FpgaProducer = ProducerOf<Fpga>;
pub const D: usize = <FpgaProducer as ProducerConfig>::DEPTH;
pub static BUF: [u32; D] = [0; D];
\end{lstlisting}

A derive on the configuration trait could emit the per-level aliases,
which would make the long spelling something nobody writes.

Wiring is not a stored back-reference. A child is handed the ports it
needs when its \code{run} is called, and that is what the inputs and
outputs of \code{Module} are for.

\begin{lstlisting}[caption={Instantiating two submodules and wiring them. Both children receive the same bus; neither holds a reference to the other.},label={lst:hier-run}]
impl<TC: TopConfig> Module<(), ()> for Top<TC> {
    async fn run(&mut self, _i: (), _o: ()) {
        join2(
            self.producer.run((), ProducerOut { out: &self.link }),
            self.consumer.run(ConsumerIn { inp: &self.link }, ()),
        ).await;
    }
}
\end{lstlisting}

\subsection{Why the borrow rule does not bite here}

Section~\ref{sec:e-units} showed that two processes of one unit cannot
both take \code{\&mut self}. Listing~\ref{lst:hier-run} appears to do
exactly that and compiles, which is worth explaining rather than leaving
as an inconsistency.

The two cases are not the same. A process is a method on the unit, so it
borrows the whole of \code{self} and two of them collide. A submodule is a
\emph{field}, and Rust permits simultaneous mutable borrows of disjoint
fields, so \code{\&mut self.producer} and \code{\&mut self.consumer} coexist
legally. The bus between them is borrowed shared by both, which is what a
wire is.

\begin{rulebox}[Where interior mutability is needed]
State that two processes of one unit both drive. Not submodules, which are
disjoint fields, and not the buses between them, which are shared by
design.
\end{rulebox}

\subsection{Arrays of instances}

An array of identical submodules is an array, and joining them is a join
over an iterator.

\begin{lstlisting}[caption={An instance array. The replication count is a Rust array length, so a config can set it.},label={lst:hier-array}]
join_all(
    self.pes.iter_mut()
        .enumerate()
        .map(|(i, pe)| pe.run(i as u32, ()))
).await;
\end{lstlisting}

TxHDL wrote this as \code{instance pe[16]} with a \code{connect} loop
beside it. Here the replication is an array length and the connection is
what each element is passed, so a config may set the count through an
associated constant and nothing else changes.
